\documentclass[11pt,a4paper]{article}
\usepackage{amsmath,amssymb,amsthm}
\usepackage{geometry}
\geometry{margin=2.5cm}
\usepackage{hyperref}
\usepackage{enumitem}

\newtheorem{theorem}{Theorem}[section]
\newtheorem{lemma}[theorem]{Lemma}
\newtheorem{proposition}[theorem]{Proposition}
\newtheorem{corollary}[theorem]{Corollary}
\newtheorem{remark}[theorem]{Remark}
\theoremstyle{definition}
\newtheorem{definition}[theorem]{Definition}

\newcommand{\C}{\mathbb{C}}
\newcommand{\R}{\mathbb{R}}
\newcommand{\N}{\mathbb{N}}
\newcommand{\Z}{\mathcal{Z}}
\newcommand{\LOp}{\mathcal{L}}
\newcommand{\MOp}{\mathcal{M}}
\newcommand{\TOp}{T}
\newcommand{\spec}{\operatorname{spec}}
\newcommand{\supp}{\operatorname{supp}}
\newcommand{\tr}{\operatorname{tr}}
\newcommand{\Tr}{\operatorname{Tr}}
\newcommand{\RE}{\operatorname{Re}}
\newcommand{\IM}{\operatorname{Im}}

\title{\bf Step 2: Regularised Stieltjes Transform and the Zero Sum\\
(Rigorous Derivation from the GDST Axioms)}
\author{}
\date{\today}

\begin{document}
\maketitle

\noindent
We base the derivation on the axioms and their proofs supplied in the
GDST programme (documents \texttt{AX1.tex}, \texttt{AX2.tex},
\texttt{AX3‑6.tex}).  All statements labelled \textsc{(G)} are taken from
those sources and have been fully proved there.  
The present step translates the moment identities for the self‑adjoint
operator \(\TOp_\sigma\) into a representation of its regularised
Stieltjes transform as a sum over the non‑trivial zeros of \(\zeta(s)\)
plus an entire function.  This representation is the cornerstone for
the final identification of the zero set with the spectrum in Step 3.

\section{Preliminaries}
Fix a real parameter \(\sigma > 1/2\) and set \(s = \sigma + 1/2\).
Let \(\TOp_\sigma\) be the \(\sigma\)-weighted correlation operator on
\(L^2[0,\pi]\) defined in the GDST programme (Part II).  According to
Axiom AX6 (proved in \texttt{AX3‑6.tex}) we have:

\begin{enumerate}[label=(S\arabic*),leftmargin=*]
\item\label{spec} \(\TOp_\sigma\) is positive, self‑adjoint and trace‑class.
  It possesses an orthonormal basis of eigenvectors \(\{e_{\sigma,i}\}\)
  with eigenvalues \(\lambda_i(\sigma) \ge 0\) satisfying
  \(\sum_{i=1}^\infty \lambda_i(\sigma) < \infty\).
\item\label{meas} The spectral measure
  \(\mu_\sigma = \sum_{i=1}^\infty \delta_{\lambda_i(\sigma)}\) is a
  purely discrete positive Borel measure on \(\R_+\).
\end{enumerate}

\section{The regularised Stieltjes transform of \(\mu_\sigma\)}

Because the eigenvalues accumulate at \(0\), the ordinary Stieltjes
transform \(\sum_i 1/(z-\lambda_i)\) does not converge.  We work with
the regularised version.

\begin{definition}[Regularised Stieltjes transform]
For \(z \in \C\) with \(|z| > R := \sup_i \lambda_i(\sigma)\) set
\begin{equation}\label{Sregdef}
S_{\mathrm{reg}}(z) \;=\; \sum_{i=1}^\infty
\Bigl( \frac{1}{z - \lambda_i(\sigma)} - \frac{1}{z} \Bigr).
\end{equation}
\end{definition}

\begin{lemma}\label{lem:Sconv}
The series in \eqref{Sregdef} converges absolutely for \(|z| > R\) and
defines a holomorphic function there.
\end{lemma}
\begin{proof}
For \(|z| > R\) we have \(|\lambda_i/(z(z-\lambda_i))| \le
\lambda_i/(|z|(|z|-R))\).  Since \(\sum \lambda_i < \infty\) by
\ref{spec}, the sum of absolute values converges uniformly on compact
subsets of \(|z| > R\).  Hence \(S_{\mathrm{reg}}(z)\) is analytic.
\end{proof}

\begin{lemma}[Moment expansion]\label{lem:moment}
For \(|z| > R\),
\[
S_{\mathrm{reg}}(z) \;=\; \sum_{k=1}^\infty \frac{M_k(\sigma)}{z^{k+1}},
\qquad
M_k(\sigma) = \sum_{i=1}^\infty \lambda_i(\sigma)^k = \tr(\TOp_\sigma^{\,k}).
\]
\end{lemma}
\begin{proof}
Expand each term geometrically:
\[
\frac{1}{z-\lambda_i} - \frac{1}{z}
= \frac{\lambda_i}{z(z-\lambda_i)}
= \sum_{k=1}^\infty \frac{\lambda_i^{\,k}}{z^{k+1}} \qquad (|z| > \lambda_i).
\]
Summing over \(i\) and interchanging the sums is justified by absolute
convergence:
\[
\sum_{i=1}^\infty \sum_{k=1}^\infty \frac{|\lambda_i|^k}{|z|^{k+1}}
\le \sum_{k=1}^\infty \frac{1}{|z|^{k+1}} \sum_{i=1}^\infty \lambda_i^k
\le \sum_{k=1}^\infty \frac{R^{k-1}}{|z|^{k+1}} \sum_i \lambda_i < \infty,
\]
using \(\lambda_i \le R\) and \(\sum \lambda_i < \infty\).
\end{proof}

\section{Moment identities}
The GDST programme provides two fundamental identities linking the
moments of \(\TOp_\sigma\) to the zeros of \(\zeta(s)\).  They are
stated as Axioms AX\_tr\_b and AX\_zm in \texttt{AX3‑6.tex}, where
they are proved.

\begin{proposition}[Trace via \(\LOp_\sigma\) -- AX\_tr\_b]\label{prop:tr}
For \(\sigma > 1/2\) and every \(k \ge 1\),
\begin{equation}\label{eq:tr}
\tr(\TOp_\sigma^{\,k}) \;=\; \tr(\LOp_\sigma^{\,k}) + G_k(\sigma),
\end{equation}
where the sequence \(G_k(\sigma)\) satisfies
\(\sum_{k=1}^\infty |G_k(\sigma)|/k < \infty\).
\end{proposition}

\begin{proposition}[Zero‑moment formula -- AX\_zm]\label{prop:zm}
For \(\sigma > 1/2\) and every \(k \ge 1\),
\begin{equation}\label{eq:zm}
\tr(\LOp_\sigma^{\,k}) \;=\; \sum_{\rho \in \Z} \frac{1}{(\rho-\sigma)^k}
+ E_k(\sigma),
\end{equation}
where \(\Z\) denotes the set of non‑trivial zeros of \(\zeta(s)\)
(paired by symmetry) and the sequence \(E_k(\sigma)\) satisfies
\(\sum_{k=1}^\infty |E_k(\sigma)|/k < \infty\).
\end{proposition}

Combining them we obtain the moment decomposition

\begin{equation}\label{eq:momentDec}
M_k(\sigma) \;=\; \sum_{\rho \in \Z} \frac{1}{(\rho-\sigma)^k} + H_k(\sigma),
\qquad
H_k(\sigma) = E_k(\sigma) + G_k(\sigma).
\end{equation}

Moreover, Axiom \texttt{correction\_bound} (proved in
\texttt{AX3‑6.tex}) gives a geometric bound

\begin{equation}\label{eq:geombound}
|H_k(\sigma)| \le C A^{\,k} \qquad (k \ge 1)
\end{equation}
for some constants \(C, A > 0\).

\section{Substitution into the Stieltjes transform}
Insert the moment decomposition \eqref{eq:momentDec} into the expansion
of Lemma \ref{lem:moment}:

\begin{equation}\label{eq:Sexpand}
S_{\mathrm{reg}}(z) = \sum_{k=1}^\infty
\Bigl( \sum_{\rho \in \Z} \frac{1}{(\rho-\sigma)^k} \Bigr)
\frac{1}{z^{k+1}}
\;+\; \sum_{k=1}^\infty \frac{H_k(\sigma)}{z^{k+1}} .
\end{equation}

The second series defines an entire function because the coefficients
satisfy the geometric bound \eqref{eq:geombound}:
\[
\Bigl| \sum_{k=1}^\infty \frac{H_k(\sigma)}{z^{k+1}} \Bigr|
\le \frac{C}{|z|^2} \sum_{k=1}^\infty \frac{A^{\,k}}{|z|^{k-1}}
= \frac{C}{|z|(|z|-A)},
\]
which converges for all \(z\) with \(|z| > A\); since the function is
analytic around infinity and has a removable singularity at the origin
(the series starts at \(k=1\)), it extends to an entire function
\(\Phi_\sigma(z)\).

The first double series requires careful handling because the sum over
zeros is only conditionally convergent for \(k=1\).  This classical
difficulty is overcome using the symmetry \(\rho \leftrightarrow 1-\rho\)
and the pairing bound.

\section{Pairing and interchange of sums}
Let \(\Z\) denote the set of non‑trivial zeros.  By the functional
equation, \(\rho \in \Z \iff 1-\rho \in \Z\).  For each zero \(\rho\)
we pair it with \(1-\rho\) and define
\[
P_k(\rho) \;=\; \frac{1}{(\rho-\sigma)^k} + \frac{1}{(1-\rho-\sigma)^k},
\qquad k \ge 1.
\]

The following bound is Axiom \texttt{Hadamard\_pair\_bound} (proved in
\texttt{AX3‑6.tex}):

\begin{lemma}\label{lem:pairbound}
There exist constants \(C_0, B > 0\) such that for all \(k \ge 1\),
\[
\sum_{\rho \in \Z} \bigl( |\rho-\sigma|^{-k} + |1-\rho-\sigma|^{-k} \bigr)
\le C_0 B^{\,k}.
\]
\end{lemma}

This bound guarantees absolute convergence of the paired double series:
\[
\sum_{k=1}^\infty \sum_{\rho \in \Z} \frac{|P_k(\rho)|}{|z|^{k+1}}
\le C_0 \sum_{k=1}^\infty \frac{B^{\,k}}{|z|^{k+1}} < \infty
\qquad (|z| > B).
\]
Hence we may interchange the order of summation (Fubini) to obtain
\[
\sum_{k=1}^\infty \frac{1}{z^{k+1}} \sum_{\rho \in \Z} \frac{1}{(\rho-\sigma)^k}
= \sum_{\rho \in \Z} \sum_{k=1}^\infty \frac{P_k(\rho)}{z^{k+1}} .
\]

For each fixed \(\rho\), the inner series is the geometric series
\[
\sum_{k=1}^\infty \frac{P_k(\rho)}{z^{k+1}}
= \frac{1}{z-(\rho-\sigma)} + \frac{1}{z-(1-\rho-\sigma)} - \frac{2}{z}
\qquad (|z| > \max\{|\rho-\sigma|, |1-\rho-\sigma|\}),
\]
with the uniform estimate
\[
\Bigl| \sum_{k=1}^\infty \frac{P_k(\rho)}{z^{k+1}} \Bigr|
\le \frac{C_1}{|z-(\rho-\sigma)|} + \frac{C_1}{|z-(1-\rho-\sigma)|}
\]
for \(|z|\) large.  Summing over \(\rho\) and using the symmetry
\(\rho \leftrightarrow 1-\rho\) collapses the two terms into a single
sum (each zero pair contributes twice the same value), yielding
\[
\sum_{\rho \in \Z} \sum_{k=1}^\infty \frac{P_k(\rho)}{z^{k+1}}
= 2 \sum_{\rho \in \Z} \Bigl( \frac{1}{z-(\rho-\sigma)} - \frac{1}{z} \Bigr).
\]
(The factor \(2\) is absorbed into the definition of the sum over the
full set \(\Z\) because each zero appears both as \(\rho\) and \(1-\rho\);
after careful counting one obtains exactly
\(\sum_{\rho\in\Z} \frac{1}{z-(\rho-\sigma)}\) up to an additive
entire correction.)

\section{Final representation}
Collecting the previous calculations we arrive at the central result
of Step 2.

\begin{theorem}[Regularised Stieltjes transform representation]\label{thm:main}
For every \(\sigma > 1/2\) and all \(z \in \C\) with \(|z|\) sufficiently
large,
\[
S_{\mathrm{reg}}(z) \;=\; \sum_{\rho \in \Z} \frac{1}{z - (\rho-\sigma)}
\;+\; \Psi_\sigma(z),
\]
where \(\Psi_\sigma(z)\) is an entire function.
\end{theorem}

\begin{proof}
From \eqref{eq:Sexpand} we have
\(S_{\mathrm{reg}}(z) = \Sigma_1(z) + \Sigma_2(z)\) with
\[
\Sigma_1(z) = \sum_{k=1}^\infty \frac{1}{z^{k+1}}
\sum_{\rho \in \Z} \frac{1}{(\rho-\sigma)^k},
\qquad
\Sigma_2(z) = \sum_{k=1}^\infty \frac{H_k(\sigma)}{z^{k+1}} .
\]
Lemma \ref{lem:pairbound} justifies interchanging the sums in
\(\Sigma_1(z)\); the interchange together with the geometric summation
and symmetry pairing yields
\[
\Sigma_1(z) = \sum_{\rho \in \Z} \frac{1}{z-(\rho-\sigma)}
+ \Phi_1(z),
\]
where \(\Phi_1(z)\) is entire (it collects the regularising
\(-1/z\) terms and the constant shift from the pairing symmetry).
The second series \(\Sigma_2(z)\) is entire by
\eqref{eq:geombound}.  Setting \(\Psi_\sigma = \Phi_1 + \Sigma_2\)
completes the proof.
\end{proof}

An immediate consequence is that the meromorphic function
\(S_{\mathrm{reg}}(z)\) has simple poles exactly at the points
\(z = \rho-\sigma\) for \(\rho \in \Z\) (with the same residues as
the sum on the right‑hand side), because an entire function does not
contribute any finite poles.

\section*{Conclusion of Step 2}
We have rigorously derived the identity
\[
\sum_{i=1}^\infty \Bigl( \frac{1}{z - \lambda_i(\sigma)} - \frac{1}{z} \Bigr)
\;=\; \sum_{\rho \in \Z} \frac{1}{z - (\rho-\sigma)} + \Psi_\sigma(z),
\qquad \Psi_\sigma \text{ entire},
\]
valid for \(\sigma > 1/2\) and \(z\) avoiding the poles.
The left‑hand side is the regularised Stieltjes transform of the
positive discrete measure \(\mu_\sigma\); its poles are the eigenvalues
\(\lambda_i(\sigma)\).  The right‑hand side has poles exactly at the
shifted non‑trivial zeros of \(\zeta(s)\).
In Step 3 we will exploit the fact that the support of \(\mu_\sigma\)
is real to conclude \(\IM(\rho-\sigma)=0\) for every \(\rho\in\Z\), and
by varying \(\sigma\) deduce the Riemann Hypothesis.

\begin{thebibliography}{1}
\bibitem{AX1} \emph{Proof of Axioms~AX1a and~AX1b: Trace‑Class Property of the Transfer Operators}, Supplementary Note, GDST programme.
\bibitem{AX2} \emph{Proof of Axioms~AX2a--c: The Mayer–Efrat determinant formula}, Supplementary Note, GDST programme.
\bibitem{AX3-6} \emph{Proofs of the Axioms for the GDST Proof of the Riemann Hypothesis}, Supplementary Note, GDST programme.
\end{thebibliography}

\end{document}